% !TEX root = ../main.tex

\chapter{Derived manifolds}
\label[chapter]{chap:From_Smooth_Algebra_To_Derived_Manifolds}

The preceding chapter established the structural theorem
\[
  \Cinfty(-)\colon
  \mathrm{Mfd}\catop
  \longrightarrow
  \Cinfty\text{-}\mathrm{Ring}^{\mathrm{fp}}.
\]
This functor is fully faithful, and it sends transverse pullback squares of
manifolds to pushout squares of $\Cinfty$-rings. Thus smooth algebra already
contains every manifold and every transverse intersection. It contains more:
the nontransverse line--parabola intersection produced the finitely presented
ring
\[
  \Cinfty(\R)/(x^2),
\]
which is not the function ring of a manifold.

These examples suggest a universal enlargement of the category of manifolds.

\begin{question}
  \label[question]{question:Universal_Finite_Limit_Completion}
  What is the universal way to enlarge manifolds so that all finite limits
  exist, while preserving transverse pullbacks and the terminal object?
\end{question}

The answer depends on the categories in which we ask the question.
Ordinary smooth rings retain nilpotent thickenings, but their pushouts
collapse the excess self-intersections seen in
\Cref{ex:Excess_Self_Intersection}. To retain that information, we pass
to $\infty$-categories, where morphisms carry identifications between them,
identifications between those identifications, and so on. Pullbacks can
then retain this higher information.

We first describe the universal construction for ordinary smooth rings
and locate its limitation in a self-intersection calculation. We then
introduce the pullback formalism in $\infty$-categories and extend smooth
algebra to this setting. The final step relates the resulting animated smooth rings
to a universal category of derived manifolds. We quote the algebraic
and geometric extension theorems, explaining their content through
presentations and examples. The next chapter will calculate the resulting
derived intersections.

\section{The ordinary finite-limit envelope}
\label[section]{sec:Ordinary_Finite_Limit_Envelope}

We begin with Euclidean spaces and freely adjoin finite limits while
preserving their finite products. The
result is the opposite category of finitely presented smooth rings:
generators specify coordinates, and relations specify equations.
Its universal property says that a finite-limit-preserving functor out
of this category is determined by an interpretation of smooth operations
in its target.

Recall that $\mathrm{Euc}$ consists of Euclidean spaces and smooth
maps. In \Cref{prop:Cinfty_Ring_Functor_Formulation}, an ordinary
$\Cinfty$-ring was a finite-product-preserving functor
$\mathrm{Euc}\to\mathrm{Set}$. The same definition works in any ordinary
category with finite products.

\begin{definition}[$\Cinfty$-ring object]
  \label[definition]{def:Internal_Cinfty_Ring}
  Let $C$ be an ordinary category with finite products. A \emph{$\Cinfty$-ring object in
  $C$} is a finite-product-preserving functor
  \[
    R\colon\mathrm{Euc}\longrightarrow C.
  \]
  We write $\operatorname{Alg}_{\Cinfty}(C)$ for the category of
  $\Cinfty$-ring objects in $C$.
\end{definition}

The object $R(\R)$ carries an operation in $C$ for every smooth map between
Euclidean spaces. Since $R(\R^n)\cong R(\R)^n$, the functor is determined by
this one object together with its smooth operations. For $C=\mathrm{Set}$ we
recover ordinary $\Cinfty$-rings.

A finite-limit-preserving functor transports such a ring object,
because it preserves finite products and the identities between the
smooth operations. We will use this to characterize an extension
from Euclidean spaces by its value on the smooth line.

To describe the universal target for these operations, set
\[
  \operatorname{Aff}_{\Cinfty}^{\mathrm{fp}}
  =
  (\Cinfty\text{-}\mathrm{Ring}^{\mathrm{fp}})\catop.
\]
An object of this opposite category is an affine smooth object described
by its function ring. Its spectrum and structure sheaf will be constructed
in \Cref{chap:Local_Derived_Geometry_Kuranishi_Charts}.

There is a functor
\[
  j\colon\mathrm{Euc}
  \longrightarrow
  \operatorname{Aff}_{\Cinfty}^{\mathrm{fp}}
\]
which sends $\R^n$ to the free ring $\Cinfty(\R^n)$, regarded as an object of
the opposite category. A smooth map $f\colon\R^n\to\R^m$ induces
\[
  f^*\colon\Cinfty(\R^m)\longrightarrow\Cinfty(\R^n),
\]
which becomes a morphism in the same direction as $f$ after taking opposites.
The functor $j$ preserves finite products.

The category of formal affine smooth objects has the following universal
property.

\begin{theorem}[Ordinary finite-limit envelope]
  \label[theorem]{thm:Ordinary_Cinfty_Finite_Limit_Envelope}
  Let $C$ be an ordinary category with finite limits. Restriction along $j$
  induces an equivalence of categories
  \[
    \Fun^{\mathrm{lex}}
    \bigl(\operatorname{Aff}_{\Cinfty}^{\mathrm{fp}},C\bigr)
    \simeq
    \operatorname{Alg}_{\Cinfty}(C),
  \]
  where the left side consists of finite-limit-preserving functors.
\end{theorem}

This is the $0$-truncated case of the finite-limit-envelope theorem for an
algebraic theory
\cite[Theorem~3.22 and Remark~3.26]{Carchedi_Steffens_Universal_Derived_Manifolds}.

The universal property has an explicit consequence: a functor preserving
finite limits must send a presentation by equations to the corresponding
zero locus in its target. A finitely presented $\Cinfty$-ring can be written
\[
  A=
  \frac{\Cinfty(\R^n)}{(f_1,\ldots,f_k)}.
\]
Equivalently, $A$ is the coequalizer
\[
  \begin{tikzcd}[column sep=large]
    \Cinfty(\R^k)
      \ar[r,shift left=.6ex,"{(f_1,\ldots,f_k)^*}"]
      \ar[r,shift right=.6ex,"0^*"']
      & \Cinfty(\R^n) \rar & A,
  \end{tikzcd}
\]
where the $i$th coordinate function on $\R^k$ is sent respectively to $f_i$
and to zero. In the opposite category, this becomes an equalizer diagram
\[
  \begin{tikzcd}[column sep=large]
    A \rar & \Cinfty(\R^n)
      \ar[r,shift left=.6ex]
      \ar[r,shift right=.6ex]
      & \Cinfty(\R^k).
  \end{tikzcd}
\]
Consequently, if $F$ preserves finite limits and
\[
  R=F\circ j\colon\mathrm{Euc}\longrightarrow C,
\]
then $F(A)$ is forced to be the internal zero locus
\[
  F(A)
  \cong
  \operatorname{Eq}
  \left(
    \begin{tikzcd}[column sep=large]
      R(\R)^n
        \ar[r,shift left=.6ex,"{(f_1,\ldots,f_k)}"]
        \ar[r,shift right=.6ex,"0"']
        & R(\R)^k
    \end{tikzcd}
  \right).
\]
Thus the chosen $\Cinfty$-ring object determines the value of $F$ on every
finite presentation.

This argument explains uniqueness, but a theorem is still required. One must
show that the construction is independent of the chosen presentation,
functorial in $A$, and compatible with natural transformations. The general
proof observes that finitely presented algebras are generated from finitely
generated free algebras by finite colimits and retracts. Taking opposites
turns those operations into finite limits and retracts. The same generation
statements underlie the animated version of the theorem. Extending from
Euclidean spaces to all manifolds requires a geometric theorem, which we
will state in \Cref{sec:Universal_Property_Derived_Manifolds}.

\section{Why ordinary intersections are not enough}
\label[section]{sec:What_Ordinary_Envelope_Forgets}

The ordinary envelope gives a formal zero locus for every smooth map,
even when the geometric zero set is singular. Quotients retain the
equations and can therefore distinguish nilpotent thickenings. They do
not, however, distinguish a closed subobject from its self-intersection.
This calculation will explain what needs to change in the universal
problem.

Let $M$ be a manifold and let
\[
  f=(f_1,\ldots,f_k)\colon M\longrightarrow\R^k
\]
be smooth. The diagram of function rings is
\[
\begin{tikzcd}
  \Cinfty(\R^k) \rar["f^*"] \dar["0^*"'] \drar[pushout]
    & \Cinfty(M) \dar \\
  \R \rar & A_f.
\end{tikzcd}
\]
By the quotient formula from \Cref{prop:Cinfty_Coproduct_Formulas},
\[
  A_f
  \cong
  \frac{\Cinfty(M)}{(f_1,\ldots,f_k)}.
\]
After taking opposites, the corresponding formal affine object is the
pullback of
\[
  M\xrightarrow{f}\R^k\xleftarrow{0}*.
\]
This construction makes sense without a spectrum or a structure sheaf.

\begin{example}[A transverse zero]
  \label[example]{ex:Ordinary_Envelope_Transverse_Zero}
  The quotient $\Cinfty(\R)/(x)\cong\R$ says that the formal zero locus
  of $x$ is the ordinary point, as required by transverse compatibility.
\end{example}

\begin{example}[A singular zero]
  \label[example]{ex:Ordinary_Envelope_Singular_Zero}
  The dual-number quotient $\Cinfty(\R)/(x^2)$ from the preceding
  chapters is the formal zero locus of $x^2$. Thus the ordinary envelope
  already retains its nilpotent thickening.
\end{example}

Now consider the self-intersection of a closed subobject. Let $A$ be a
$\Cinfty$-ring, let $I\subseteq A$ be an ideal, and put
\[
  B=A/I.
\]
The self-intersection of the corresponding formal closed subobject is computed
by
\[
  B\amalg_A B.
\]

\begin{proposition}[Ordinary self-intersections]
  \label[proposition]{prop:Ordinary_Self_Intersections_Collapse}
  There is a natural isomorphism
  \[
    (A/I)\amalg_A(A/I)
    \cong
    A/I.
  \]
\end{proposition}

\begin{proof}
  Apply the quotient pushout formula:
  \[
    (A/I)\amalg_A(A/I)
    \cong
    A/(I+I)
    =A/I.
  \]
\end{proof}

For example, let $A=\Cinfty(\R^2)$ and $I=(y)$. Then $B$ represents the
$x$-axis. Its ordinary self-intersection inside the plane is again represented
by $B$. The quotient remembers that $y$ vanishes, but it has no record that
the relation $y=0$ was imposed twice.

The collapse reflects how pullbacks work in an ordinary category.
An ordinary pullback is characterized by an isomorphism of sets
\[
  \Hom_C(T,X\times_ZY)
  \cong
  \Hom_C(T,X)
  \times_{\Hom_C(T,Z)}
  \Hom_C(T,Y).
\]
An element on the right is a pair of maps whose composites to $Z$ are equal.
In a set, equality is a property. Once the equality holds, there is no further
datum recording how it holds. In particular, imposing the same relation twice
cannot create new equality data.

\section{Pullbacks with higher identifications}
\label[section]{sec:Equality_As_Structure}

We want a pullback to record how two composites are identified. For this,
we first introduce animae, whose elements admit identifications and
higher identifications. The based loops in a circle will illustrate
what a self-pullback can retain. We then formulate pullbacks in
$\infty$-categories before returning to smooth algebra.

\begin{definition}[Anima, informal description]
\label[definition]{def:Anima_Conceptual}
An \emph{anima} behaves like a set in which the identifications between
any two elements themselves form an anima. Identifications can be composed
and inverted, with the usual identities holding through coherent higher
identifications.
\end{definition}

For the purposes of this seminar, we take the notion of an anima as a
foundational concept, much as we take the notion of a set in ordinary
mathematics. It is possible to construct explicit models using, for example,
simplicial sets. Such constructions require substantial setup and do not
by themselves convey the intended intuition. We will instead work with
the heuristic picture above: an anima behaves like a set, except that its
identifications themselves form animae.

The following examples give different meanings to these identifications.
\begin{itemize}
  \item Every set $X$ determines an anima. For $x,y\in X$, the anima of
    identifications between $x$ and $y$ is a point if $x=y$, and empty
    otherwise. These are the $0$-truncated animae.
  \item Every ordinary category $C$ determines an anima $C^{\simeq}$,
    its \emph{groupoid core}. Its points are the objects of $C$, and its
    identifications are isomorphisms in $C$. For example, there is an
    anima of manifolds whose identifications are diffeomorphisms.
  \item Every topological space $X$ determines an anima $\Pi_\infty(X)$.
    Its points are points of $X$, and its identifications are paths in $X$.
    Here the anima of identifications can itself be interesting:
    identifications between paths are homotopies relative to their
    endpoints, and continuing in this way gives iterated homotopies.
\end{itemize}

An ordinary category can be described by a set of objects and a set of
morphisms between each pair of objects. In an $\infty$-category $C$, we
instead have an anima of objects and, for every two objects $X,Y$, an anima
\[
  \Hom_C(X,Y)
\]
of morphisms. An identification in $\Hom_C(X,Y)$ is called a
\emph{homotopy} between morphisms. Composition is associative and unital
with coherent higher identifications. Universal properties that are
ordinarily phrased using equalities between morphisms are now phrased
using homotopies between morphisms, together with their higher coherence.

Ordinary categories are the special case in which every mapping anima is
a set, so there are no nontrivial homotopies between morphisms. When an
ordinary category $C$ is viewed in this way, its anima of objects is its
core $C^{\simeq}$; isomorphisms between objects are still retained.

We can now write $\An$ for the $\infty$-category of animae. We will use
products, pullbacks, functor categories, and equivalences in this setting
without choosing a model.

The effect of retaining identifications is already visible in a
self-pullback of animae. Let $X$ be an anima and let $x\colon *\to X$
be a point. Form the pullback
\[
\begin{tikzcd}
  *\times_X* \rar \dar \drar[pullback]
    & * \dar["x"] \\
  * \rar["x"'] & X.
\end{tikzcd}
\]
An element of this pullback consists of the two chosen points together with
an identification between their images in $X$. Both images are $x$.
Writing $X(x,x)$ for the anima of identifications of $x$ with itself,
including all higher identifications, we obtain an isomorphism
\[
  *\times_X*
  \simeq
  X(x,x).
\]

In sets, the right side is a singleton. In animae, it can be highly
nontrivial. If
\[
  X=\Pi_\infty(S^1)
\]
is the underlying anima of the circle with base point $x$, an identification
of $x$ with itself is represented by a continuous path
\[
  \gamma\colon[0,1]\longrightarrow S^1,
  \qquad \gamma(0)=\gamma(1)=x.
\]
The endpoint condition is a strict equality of points, but the datum in the
pullback is the entire path $\gamma$. It may travel around the circle any
number of times. Paths between these choices are homotopies of loops that
keep the base point fixed. Thus
\[
  X(x,x)
  \simeq
  \Omega\Pi_\infty(S^1)
  \simeq
  \Z.
\]
The integers record the winding numbers of loops. Each winding-number
component is contractible, which gives the last isomorphism with the
discrete anima $\Z$.

The loop anima is a prototype for the information a homotopy pullback
can retain. A geometric derived self-intersection uses the same
pullback formalism in a different category.

The same description of pullbacks applies in any $\infty$-category.
Let
\[
  X\xrightarrow{f}Z\xleftarrow{g}Y
\]
be a diagram in an $\infty$-category $C$. Its pullback is characterized by
natural isomorphisms of animae
\[
  \Hom_C(T,X\times_ZY)
  \simeq
  \Hom_C(T,X)
  \times_{\Hom_C(T,Z)}
  \Hom_C(T,Y)
\]
for all $T\in C$.

Informally, a generalized $T$-point consists of
\[
  a\colon T\to X,
  \qquad
  b\colon T\to Y,
  \qquad
  \alpha\colon f\circ a\cong g\circ b.
\]
The equality $\alpha$ is part of the datum. Equalities between choices of
$\alpha$, and all subsequent coherence, are contained in the pullback of hom
animae.

This is the only pullback notion needed in the sequel. In older sources the
phrase \emph{homotopy pullback} is often used to distinguish a construction in
a chosen presentation from an ordinary strict pullback. We will work
synthetically and simply say \emph{pullback} in an $\infty$-category.

The passage from hom sets to hom animae now gives a precise reason to reconsider
\Cref{question:Universal_Finite_Limit_Completion}. We should ask for a
universal finite-limit envelope among $\infty$-categories, not merely among
ordinary categories.

\section{Animated smooth rings and finite limits}
\label[section]{sec:Animated_Finite_Limit_Envelope}

We now apply the preceding pullback formalism to smooth algebra. The
operations are still indexed by smooth maps between Euclidean spaces,
but their values are animae. The finite-limit-envelope theorem has a
corresponding animated form, with compactness replacing ordinary finite
presentation and with retracts included explicitly.

The definition in \Cref{def:Internal_Cinfty_Ring} extends to an $\infty$-category $C$
with finite products: a smooth ring object is a finite-product-preserving
functor $\mathrm{Euc}\to C$, with $\mathrm{Euc}$ regarded as an
$\infty$-category. We retain the notation $\operatorname{Alg}_{\Cinfty}(C)$.
Taking $C=\An$ gives the following definition.

\begin{definition}[Animated $\Cinfty$-ring]
  \label[definition]{def:Animated_Cinfty_Ring}
  An \emph{animated $\Cinfty$-ring} is a finite-product-preserving functor
  \[
    A\colon\mathrm{Euc}\longrightarrow\An.
  \]
  We denote their $\infty$-category by
  \[
    \operatorname{Alg}_{\Cinfty}(\An)
    =
    \Fun^\times(\mathrm{Euc},\An).
  \]
\end{definition}

An ordinary $\Cinfty$-ring is the special case in which every value $A(\R^n)$
is a set, regarded as a $0$-truncated anima. Animated rings allow the smooth
operations themselves to interact with higher equality data.

For ordinary smooth rings, \Cref{prop:Filtered_Colimit_Criterion} expresses
finite presentation by preservation of filtered colimits. We use the same
criterion here, applied to the whole mapping anima.

\begin{definition}[Homotopical finite presentation]
  \label[definition]{def:Homotopically_Finitely_Presented_Cinfty_Ring}
  An animated $\Cinfty$-ring $A$ is \emph{homotopically finitely presented} if
  it is compact in $\operatorname{Alg}_{\Cinfty}(\An)$: the functor
  \[
    \Hom_{\operatorname{Alg}_{\Cinfty}(\An)}(A,-)
    \colon
    \operatorname{Alg}_{\Cinfty}(\An)
    \longrightarrow
    \An
  \]
  preserves filtered colimits.
\end{definition}

We use the subscript $\mathrm{fp}$ for the full subcategory of such objects.
The word \emph{homotopically} matters: finite presentation is now a statement
about the entire hom anima, not only its set of connected components.

The next theorem identifies the universal target for Euclidean smooth
operations when finite limits and retracts are required. It applies to
every algebraic theory; we state only its $\Cinfty$ instance.

\begin{theorem}[Animated finite-limit envelope]
  \label[theorem]{thm:Animated_Cinfty_Finite_Limit_Envelope}
  Let $C$ be an idempotent-complete $\infty$-category with finite limits.
  Restriction along the finitely generated free animated $\Cinfty$-rings
  induces an equivalence of $\infty$-categories
  \[
    \Fun^{\mathrm{lex}}
    \left(
      \operatorname{Alg}_{\Cinfty}(\An)_{\mathrm{fp}}\catop,
      C
    \right)
    \simeq
    \operatorname{Alg}_{\Cinfty}(C).
  \]
\end{theorem}

This is \citet[Theorem~3.22]{Carchedi_Steffens_Universal_Derived_Manifolds}.
The ordinary \Cref{thm:Ordinary_Cinfty_Finite_Limit_Envelope} is its
$0$-truncated counterpart.

Recall that an $\infty$-category is idempotent-complete if every retract
diagram which exists formally is represented by an actual object. Ordinary
categories with finite limits are automatically idempotent-complete. This is
not automatic for $\infty$-categories, which is why the condition appears in
the theorem.

We sketch the ingredients of the quoted theorem. The starting point is
compact generation by free algebras.

\begin{enumerate}
  \item The free $\Cinfty$-ring on $n$ generators is compact. Mapping out of
    it is evaluation on an $n$-tuple, and evaluation commutes with filtered
    colimits.
  \item Every animated $\Cinfty$-ring is assembled as a colimit of free
    rings. This is the homotopical form of presenting an algebra by generators
    and relations.
  \item A compact animated $\Cinfty$-ring is a retract of a finite colimit of
    finitely generated free rings.
  \item After taking opposites, the preceding statement says that every object
    of
    $\operatorname{Alg}_{\Cinfty}(\An)_{\mathrm{fp}}\catop$
    is a retract of a finite limit of Euclidean objects. This describes the values forced on any finite-limit-preserving
    extension from $\mathrm{Euc}$.
\end{enumerate}

The complete argument constructs the extension by right Kan extension and
checks that it takes values in the chosen target $C$. The compactness and
idempotent-completeness hypotheses are precisely what make that last step
work. Finite colimits and retracts of free algebras explain the form of the
universal property; the Kan-extension argument supplies existence and
compatibility with all presentations.

Replacing sets by animae changes both the limits in the universal
problem and the finiteness condition on its algebraic objects.
The new pullbacks retain homotopies and their coherence. Compactness
must likewise control the whole mapping anima.

\begin{warning}
  \label[warning]{warn:Ordinary_Fp_Not_Necessarily_Animated_Fp}
  An ordinary finitely presented $\Cinfty$-ring need not remain compact when it
  is regarded as a discrete animated $\Cinfty$-ring. Thus the ordinary
  envelope is not simply a full subcategory of the animated envelope obtained
  by declaring every ordinary object discrete.
\end{warning}

The analogous warning for an arbitrary algebraic theory is recorded in
\citet[Warning~3.28]{Carchedi_Steffens_Universal_Derived_Manifolds}.
The truncation functor $\pi_0$ does send homotopically finitely presented
animated rings to finitely presented ordinary rings, but the reverse passage
can require retaining higher relation data.

In particular, the animated pushout associated with a self-intersection
can retain higher homotopy groups beyond its ordinary quotient.

\section{The universal property of derived manifolds}
\label[section]{sec:Universal_Property_Derived_Manifolds}

The animated envelope was constructed from Euclidean spaces. To answer
the chapter's question, we must relate its universal property to functors
on all manifolds that preserve transverse pullbacks. A geometric extension
theorem provides this link. We first state that theorem, then define the
universal category of derived manifolds and identify it with the opposite
category of compact animated smooth rings.

\Cref{chap:Manifolds_As_Cinfty_Rings} showed that the function-ring functor
preserves transverse pullbacks. We need a stronger statement: any
interpretation of Euclidean smooth operations extends uniquely to
manifolds while preserving transverse pullbacks and the terminal object.
We quote this geometric extension theorem.

Let
\[
  q\colon\mathrm{Euc}\longrightarrow\mathrm{Mfd}
\]
be the inclusion. For every
idempotent-complete $\infty$-category $C$ with finite limits, restriction along
$q$ induces an equivalence
\[
  \Fun^{\pitchfork}(\mathrm{Mfd},C)
  \simeq
  \Fun^\times(\mathrm{Euc},C)
  =
  \operatorname{Alg}_{\Cinfty}(C).
\]
Here $\Fun^\pitchfork$ denotes the functors preserving transverse
pullbacks and the terminal object. This is
\citet[Lemma~5.1 and Corollary~5.2]{Carchedi_Steffens_Universal_Derived_Manifolds}.

The geometric idea comes from the presentations in
\Cref{chap:Manifolds_As_Cinfty_Rings}. An open Euclidean subset is a
transverse zero locus of $ya(x)-1$, and every manifold is a retract of
an open tubular neighborhood in Euclidean space. A product-preserving
functor on $\mathrm{Euc}$ can therefore assign an object to each such
presentation using finite limits and retracts. This explains the candidate
extension. Proving its independence of choices, functoriality, and
preservation of transverse pullbacks is the substantive part of the
theorem; the proof uses right Kan extension.

The earlier transverse-pullback theorem concerned the actual function-ring
functor. This extension theorem applies to any internal $\Cinfty$-ring:
it determines a unique transverse-pullback-preserving functor on manifolds.

For an ordinary category $C$ with finite limits, combining this extension
theorem with \Cref{thm:Ordinary_Cinfty_Finite_Limit_Envelope} gives
\[
  \Fun^{\mathrm{lex}}(\operatorname{Aff}_{\Cinfty}^{\mathrm{fp}},C)
  \simeq\Fun^\pitchfork(\mathrm{Mfd},C).
\]
Thus ordinary finitely presented smooth rings solve the ordinary version
of the chapter's universal problem. We now formulate the version for
$\infty$-categories.

\begin{definition}[Universal $\infty$-category of derived manifolds]
  \label[definition]{def:Universal_Infinity_Category_Derived_Manifolds}
  The $\infty$-category $\mathrm{DMfd}$ of \emph{derived manifolds} is equipped
  with a functor
  \[
    i\colon\mathrm{Mfd}\longrightarrow\mathrm{DMfd}
  \]
  and is characterized by the following properties.
  \begin{enumerate}
    \item The $\infty$-category $\mathrm{DMfd}$ has finite limits and is
      idempotent-complete.
    \item The functor $i$ preserves transverse pullbacks and the terminal
      object.
    \item For every idempotent-complete $\infty$-category $C$ with finite
      limits, restriction along $i$ induces an equivalence
      \[
        \Fun^{\mathrm{lex}}(\mathrm{DMfd},C)
        \simeq
        \Fun^{\pitchfork}(\mathrm{Mfd},C),
      \]
      where the right side consists of functors preserving transverse
      pullbacks and the terminal object.
  \end{enumerate}
\end{definition}

This is \citet[Definition~2.1]{Carchedi_Steffens_Universal_Derived_Manifolds}.
The universal property determines $\mathrm{DMfd}$ uniquely.
It fixes the existing transverse pullbacks, adjoins all finite
homotopy limits, and includes their retracts.

The geometric extension theorem turns the defining universal property
under manifolds into a universal property in terms of smooth ring objects.
Comparing this with the animated envelope will both establish existence
and give an algebraic description of $\mathrm{DMfd}$.

\begin{theorem}[Universal property through $\Cinfty$-rings]
  \label[theorem]{thm:Universal_Property_Derived_Manifolds_Cinfty_Rings}
  For every idempotent-complete $\infty$-category $C$ with finite limits,
  there is a natural equivalence
  \[
    \Fun^{\mathrm{lex}}(\mathrm{DMfd},C)
    \simeq
    \operatorname{Alg}_{\Cinfty}(C).
  \]
\end{theorem}

\begin{proof}
  By the universal property under manifolds,
  \[
    \Fun^{\mathrm{lex}}(\mathrm{DMfd},C)
    \simeq
    \Fun^{\pitchfork}(\mathrm{Mfd},C).
  \]
  Restriction to Euclidean spaces gives
  \[
    \Fun^{\pitchfork}(\mathrm{Mfd},C)
    \simeq
    \Fun^\times(\mathrm{Euc},C)
    =
    \operatorname{Alg}_{\Cinfty}(C).
  \]
  Composing these equivalences proves the claim.
\end{proof}

This is \citet[Theorem~5.3]{Carchedi_Steffens_Universal_Derived_Manifolds}.

The animated finite-limit envelope has exactly the same universal property.
We therefore obtain the promised algebraic presentation.

\begin{corollary}[Derived manifolds as formal animated affines]
  \label[corollary]{cor:Derived_Manifolds_Animated_Cinfty_Rings}
  There is an equivalence of $\infty$-categories
  \[
    \mathrm{DMfd}
    \simeq
    \operatorname{Alg}_{\Cinfty}(\An)_{\mathrm{fp}}\catop.
  \]
\end{corollary}

\begin{proof}
  By
  \Cref{thm:Animated_Cinfty_Finite_Limit_Envelope,thm:Universal_Property_Derived_Manifolds_Cinfty_Rings},
  both sides are idempotent-complete $\infty$-categories with finite limits
  which universally carry a $\Cinfty$-ring object. Universal objects are
  unique, so the two $\infty$-categories are equivalent.
\end{proof}

This is \citet[Corollary~5.4]{Carchedi_Steffens_Universal_Derived_Manifolds}.
This identifies the universal category with an algebraic construction.
In particular, manifolds embed fully faithfully, transverse pullbacks
are preserved, and all finite pullbacks exist.

For a smooth map $f\colon U\to\R^r$, the function ring of its derived
zero locus is consequently the pushout
\[
 \Cinfty(U)\otimes^{\mathbb L}_{\infty,\Cinfty(\R^r)}\R
\]
in animated smooth rings. \Cref{chap:Derived_Zero_Loci_And_Tangent_Complexes} computes it by a differential
graded resolution. The category of derived schemes introduced in
\Cref{chap:Local_Derived_Geometry_Kuranishi_Charts} allows these affine objects to be described locally and glued.

\section{Further details: retracts and coherence}
\label[section]{sec:From_Smooth_Algebra_Summary_Supplementary_Materials}

The universal properties above require idempotent-complete targets.
We explain more precisely what this means in an $\infty$-category and
why finite limits alone do not supply it. The exercises revisit this
point and the self-pullback example.

Informally, an idempotent looks like a morphism $e\colon X\to X$ together with
a homotopy from $e\circ e$ to $e$. Formally, an idempotent in an
$\infty$-category is a coherent idempotent diagram: the displayed homotopy is
part of the data, together with its higher compatibilities. Splitting it means
finding an object $Y$ and maps
\[
  Y\xrightarrow{i}X\xrightarrow{r}Y
\]
together with coherent identifications of $r\circ i$ with
$\operatorname{id}_Y$ and of $i\circ r$ with the given idempotent. In an
ordinary category with finite limits, an idempotent $e\colon X\to X$
splits by taking the equalizer $i\colon Y\to X$ of $e$ and the identity.
The map $e$ factors as $ir$, and $ri=\operatorname{id}_Y$ follows
from the equalizer property. For $\infty$-categories
this is not forced by finite limits alone. The compact algebraic objects are
closed under retracts, so their opposite category is naturally
idempotent-complete. The universal property of derived manifolds must record
this closure explicitly.

\begin{exercise}
Prove the equalizer construction of a splitting just described.
Compare it with the explicit finite presentation for a retract ring
in \Cref{chap:Manifolds_As_Cinfty_Rings}.
\end{exercise}

\begin{exercise}
Use the universal property of a homotopy pullback to identify
$*\times_{BG}*$ with a discrete group $G$, where $BG$ is its
classifying anima and both maps select the base point.
Compare this with the self-pullback over a set.
\end{exercise}

The primary source is
\cite[Definition~2.1, Theorem~3.22, Remark~3.26,
and Section~5]{Carchedi_Steffens_Universal_Derived_Manifolds}.
It treats both universal properties, their geometric comparison, and
the relation with Spivak's construction.
